\documentclass[12pt, twoside]{article}
\input{/Users/chris/github/course-files/docs/ib/preamble}

\fancyhead[LE]{\thepage}
\fancyhead[RO]{Markscheme}
\fancyhead[LO]{La Scuola d'Italia / Dr.\ Huson / IB Math HL \\* 21 May 2026}

\begin{document}
\subsubsection*{5.1 Quiz: Maclaurin Series --- Markscheme}

\begin{enumerate}[itemsep=0.3cm]

%--- Problem 1: Maclaurin of (1-x)^-4 with depreciation application [7 marks] ---
%--- Source: AA HL 2025 Nov P1 TZ1 Q05 ---
\item[1.] \makebox[\linewidth][s]{[Maximum mark: 7]\hfill {\small AA HL 2025 Nov P1 TZ1 Q05}}

The first four terms of the Maclaurin series of $(1 - x)^{-4}$ are
$1 + ax + bx^2 + 20x^3$.

\textbf{(a)(i)} Show that $a = 4$. \hfill [2]

\textbf{Solution.} Let $f(x) = (1 - x)^{-4}$. Then
\[
f'(x) = 4(1 - x)^{-5} \hfill \textbf{A1}
\]
The coefficient of $x$ in the Maclaurin series is $f'(0) = 4(1)^{-5} = 4$,
so $a = 4$. \hfill \textbf{M1 AG}

\medskip

\textbf{(a)(ii)} Find the value of $b$. \hfill [3]

\textbf{Solution.}
\[
f''(x) = 20(1 - x)^{-6} \hfill \textbf{A1}
\]
Coefficient of $x^2$ is $\dfrac{f''(0)}{2!} = \dfrac{20}{2} = 10$, so
\[
b = 10. \hfill \textbf{A1\,A1}
\]
\textit{(Alternative: binomial expansion with negative integer power
gives $\binom{-4}{2}(-x)^2 = \dfrac{(-4)(-5)}{2}x^2 = 10x^2$.)}

\medskip

\textbf{(b)} Estimate the value of the car four years ago. \hfill [2]

\textbf{Solution.} Each year the value is multiplied by $0.9$, so the
value four years ago, $V_0$, satisfies $V_0 (0.9)^4 = 1000$, i.e.\
$V_0 = 1000\,(0.9)^{-4} = 1000\,(1 - 0.1)^{-4}$.

Using the expansion from part (a) with $x = 0.1$:
\[
(1 - 0.1)^{-4} \approx 1 + 4(0.1) + 10(0.1)^2 + 20(0.1)^3
= 1 + 0.4 + 0.1 + 0.02 = 1.52. \hfill \textbf{M1}
\]
\[
V_0 \approx \$1520. \hfill \textbf{A1}
\]

\bigskip
\hrule
\bigskip

%--- Problem 2: Maclaurin of sin(x^2) and product manipulation [7 marks] ---
%--- Source: AA HL 2024 May P1 TZ1 Q08 ---
\item[2.] \makebox[\linewidth][s]{[Maximum mark: 7]\hfill {\small AA HL 2024 May P1 TZ1 Q08}}

\textbf{(a)(i)} First two non-zero terms of $\sin\!\left(x^2\right)$.
\hfill [2]

\textbf{Solution.} Start from $\sin u = u - \dfrac{u^3}{3!} + \cdots$ and
substitute $u = x^2$:
\[
\sin\!\left(x^2\right) = x^2 - \frac{x^6}{6} + \cdots
\hfill \textbf{A1\,A1}
\]

\medskip

\textbf{(a)(ii)} First two non-zero terms of $\sin^2\!\left(x^2\right)$.
\hfill [3]

\textbf{Method 1 (square the series).} \hfill \textbf{M1}
\[
\left(\sin x^2\right)^2 = \left(x^2 - \frac{x^6}{6} + \cdots\right)^2
= x^4 - 2 \cdot x^2 \cdot \frac{x^6}{6} + \cdots
= x^4 - \frac{x^8}{3} + \cdots \hfill \textbf{A1\,A1}
\]

\textbf{Method 2 (double-angle identity).} Use
$\sin^2 u = \dfrac{1 - \cos 2u}{2}$ with $u = x^2$:
\[
\sin^2\!\left(x^2\right) = \frac{1}{2}\!\left(1 - \cos 2x^2\right)
= \frac{1}{2}\!\left[\frac{(2x^2)^2}{2!} - \frac{(2x^2)^4}{4!} + \cdots\right]
= x^4 - \frac{x^8}{3} + \cdots
\]

\medskip

\textbf{(b)} First two non-zero terms of
$4x \sin\!\left(x^2\right) \cos\!\left(x^2\right)$. \hfill [2]

\textbf{Method 1 (derivative recognition).} Note that
$4x \sin(x^2) \cos(x^2) = \dfrac{d}{dx}\!\left[\sin^2(x^2)\right]$.
\hfill \textbf{M1}

Differentiating the series from (a)(ii):
\[
\frac{d}{dx}\!\left(x^4 - \frac{x^8}{3} + \cdots\right)
= 4x^3 - \frac{8x^7}{3} + \cdots \hfill \textbf{A1}
\]

\textbf{Method 2 (double-angle).} Use
$4x \sin(x^2)\cos(x^2) = 2x \sin(2x^2)$, then
\[
2x\!\left[2x^2 - \frac{(2x^2)^3}{6} + \cdots\right]
= 4x^3 - \frac{8x^7}{3} + \cdots
\]

\bigskip
\hrule
\bigskip

%--- Problem 3: Induction + Maclaurin + limit [14 marks] ---
%--- Source: AA HL 2021 Nov P1 TZ0 Q11 ---
\item[3.] \makebox[\linewidth][s]{[Maximum mark: 14]\hfill {\small AA HL 2021 Nov P1 TZ0 Q11}}

\textbf{(a)} Prove by induction that
$\dfrac{d^n}{dx^n}(x^2 e^x) = [x^2 + 2nx + n(n-1)] e^x$
for $n \in \mathbb{Z}^{+}$. \hfill [7]

\textbf{Solution.}

\textit{Base case $n = 1$:}
\[
\text{LHS} = \frac{d}{dx}(x^2 e^x) = 2x e^x + x^2 e^x
= (x^2 + 2x) e^x. \hfill \textbf{A1}
\]
\[
\text{RHS} = [x^2 + 2(1)x + 1(0)] e^x = (x^2 + 2x) e^x. \hfill \textbf{A1}
\]
So the statement holds for $n = 1$.

\textit{Inductive step.} Assume the statement holds for $n = k$, i.e.\
\[
\frac{d^k}{dx^k}(x^2 e^x) = [x^2 + 2kx + k(k-1)] e^x. \hfill \textbf{M1}
\]

Differentiating both sides: \hfill \textbf{M1}
\begin{align*}
\frac{d^{k+1}}{dx^{k+1}}(x^2 e^x)
&= \frac{d}{dx}\!\left\{[x^2 + 2kx + k(k-1)] e^x\right\} \\
&= (2x + 2k) e^x + [x^2 + 2kx + k(k-1)] e^x \hfill \textbf{A1} \\
&= [x^2 + (2k + 2)x + (k^2 + k)] e^x \\
&= [x^2 + 2(k+1)x + (k+1)k] e^x. \hfill \textbf{A1}
\end{align*}

This is the statement with $n = k + 1$. Since the statement holds for
$n = 1$ and (statement at $n = k$) $\Rightarrow$ (statement at $n = k+1$),
by mathematical induction it holds for all $n \in \mathbb{Z}^{+}$.
\hfill \textbf{R1}

\medskip

\textbf{(b)} Maclaurin series of $f(x) = x^2 e^x$ up to $x^4$. \hfill [3]

\textbf{Method 1.} Using the formula from part (a) at $x = 0$:
$\dfrac{d^n}{dx^n}(x^2 e^x)\Big|_{x=0} = n(n - 1)$, so
\[
f(0) = 0,\quad f'(0) = 0,\quad f''(0) = 2,\quad
f'''(0) = 6,\quad f^{(4)}(0) = 12. \hfill \textbf{M1}
\]
Therefore
\[
f(x) = 0 + 0 + \frac{2}{2!}x^2 + \frac{6}{3!}x^3 + \frac{12}{4!}x^4 + \cdots
= x^2 + x^3 + \tfrac{1}{2}x^4 + \cdots \hfill \textbf{M1 A1}
\]

\textbf{Method 2.} Multiply the known series for $e^x$ by $x^2$:
\[
x^2\!\left(1 + x + \tfrac{x^2}{2!} + \cdots\right)
= x^2 + x^3 + \tfrac{1}{2}x^4 + \cdots
\]

\medskip

\textbf{(c)} Evaluate
$\displaystyle\lim_{x \to 0}\frac{(x^2 e^x - x^2)^3}{x^9}$. \hfill [4]

\textbf{Solution.} From (b), $x^2 e^x - x^2 \approx x^3 + \tfrac{1}{2}x^4 + \cdots$.
\hfill \textbf{M1}

So
\[
(x^2 e^x - x^2)^3 \approx \left(x^3 + \tfrac{1}{2}x^4 + \cdots\right)^3
= x^9\!\left(1 + \tfrac{1}{2}x + \cdots\right)^3, \hfill \textbf{A1}
\]
\[
\frac{(x^2 e^x - x^2)^3}{x^9} \approx \left(1 + \tfrac{1}{2}x + \cdots\right)^3
\to 1 \quad\text{as } x \to 0. \hfill \textbf{A1\,A1}
\]
\[
\boxed{\lim_{x \to 0}\frac{(x^2 e^x - x^2)^3}{x^9} = 1.}
\]

\end{enumerate}
\end{document}
